\sectionkicker{Source-backed technical notes}
\subsection*{Agent Safety \& Security — 記憶・retrieval・harnessが新しい攻撃面になる: Technical Notes}
本欄は記事本文で圧縮した一次資料上の情報を、比較・再検証しやすい形で整理したものである。時系列、確認済みの主張、留意点、一次資料URLを掲載する。
\medskip
\subsection*{Theme at a glance}
\addcontentsline{toc}{subsection}{Theme at a glance}
\begin{center}
\footnotesize
\begin{tabularx}{\linewidth}{@{}p{0.20\linewidth}p{0.10\linewidth}p{0.14\linewidth}X@{}}
\toprule
資料 & 位置づけ & 種別 & 時系列 \\
\midrule
From Prompt Injection to Persistent Control: Defending Agentic Harness Against Trojan Backdoors & 主要資料 & 論文 & 2026-05-29 (論文公開) \\
Hijacking Agent Memory: Stealthy Trojan Attacks Through Conversational Interaction & 主要資料 & 論文 & 2026-05-28 (論文公開) \\
Relevance as a Vulnerability: How Web Retrieval Degrades Safety Alignment in LLM Agents & 補足資料 & 論文 & 2026-05-28 (論文公開) \\
Send a SCOUT First: Pre-hoc Reasoning for Adaptive Detector Allocation in Prompt-Injection Defense & 補足資料 & 論文 & 2026-05-29 (論文公開) \\
\bottomrule
\end{tabularx}
\end{center}
\normalsize

\begin{technicalnote}{From Prompt Injection to Persistent Control: Defending Agentic Harness Against Trojan Backdoors}{主要資料}
\begingroup
% reader-facing Technical Notes break policy
\widowpenalty=10000
\clubpenalty=10000
\displaywidowpenalty=10000
\begin{tabularx}{\linewidth}{@{}>{\bfseries}p{0.22\linewidth}X@{}}
組織 & - \\
種別 & 論文 \\
時系列 & 2026-05-29 (論文公開) \\
\end{tabularx}
\smallskip
{\bfseries 一次資料から整理したtechnical points}
\begin{itemize}[leftmargin=1.5em,itemsep=0.35em]
\item \textbf{一次情報で確認できる事実}: 保存済み一次資料では「From Prompt Injection to Persistent Control: Defending Agentic Harness Against Trojan Backdoors」が2026-05-29に確認できる。
\end{itemize}
\begin{minipage}{\linewidth}
% reader-facing Technical Notes coherent tail group
\begin{itemize}[leftmargin=1.5em,itemsep=0.35em]
\item \textbf{Author claim}: ClawTrojanはlocal agentic harnessに対するmulti-step persistent prompt-injection attackをモデル化し、DASGuardはworkspace内のuntrustedなcontrol-like contentを追跡・除去する方式として提示されている。
\end{itemize}
{\bfseries 読む際の境界}
\begin{itemize}[leftmargin=1.5em,itemsep=0.35em]
\item \textbf{分析上の留意点}: arXivの原メタデータ／abstractは一次論文資料として扱うが、実験結果と一般化可能性はAuthor claimであり、本号では独立再現していない。
\end{itemize}
\begin{samepage}
{\bfseries 一次資料}
\begingroup\sloppy
\begin{itemize}[leftmargin=1.5em,itemsep=0.25em]
\item {\scriptsize\url{http://arxiv.org/abs/2605.31042v1}}
\end{itemize}
\endgroup
\end{samepage}
\end{minipage}
\endgroup
\end{technicalnote}
\medskip

\begin{technicalnote}{Hijacking Agent Memory: Stealthy Trojan Attacks Through Conversational Interaction}{主要資料}
\begingroup
% reader-facing Technical Notes break policy
\widowpenalty=10000
\clubpenalty=10000
\displaywidowpenalty=10000
\begin{tabularx}{\linewidth}{@{}>{\bfseries}p{0.22\linewidth}X@{}}
組織 & - \\
種別 & 論文 \\
時系列 & 2026-05-28 (論文公開) \\
\end{tabularx}
\smallskip
{\bfseries 一次資料から整理したtechnical points}
\begin{itemize}[leftmargin=1.5em,itemsep=0.35em]
\item \textbf{一次情報で確認できる事実}: 保存済み一次資料では「Hijacking Agent Memory: Stealthy Trojan Attacks Through Conversational Interaction」が2026-05-28に確認できる。
\item \textbf{Author claim}: MemPoisonは、conversational interactionを通じてselective long-term memory pipelineを攻撃し、trigger/payload binding、entity masquerading、embedding optimizationを用いる方式を提示している。
\end{itemize}
{\bfseries 読む際の境界}
\begin{itemize}[leftmargin=1.5em,itemsep=0.35em]
\item \textbf{分析上の留意点}: arXivの原メタデータ／abstractは一次論文資料として扱うが、実験結果と一般化可能性はAuthor claimであり、本号では独立再現していない。
\end{itemize}
\begin{samepage}
{\bfseries 一次資料}
\begingroup\sloppy
\begin{itemize}[leftmargin=1.5em,itemsep=0.25em]
\item {\scriptsize\url{http://arxiv.org/abs/2605.29960v1}}
\end{itemize}
\endgroup
\end{samepage}
\endgroup
\end{technicalnote}
\medskip

\begin{technicalnote}{Relevance as a Vulnerability: How Web Retrieval Degrades Safety Alignment in LLM Agents}{補足資料}
\begingroup
% reader-facing Technical Notes break policy
\widowpenalty=10000
\clubpenalty=10000
\displaywidowpenalty=10000
\begin{tabularx}{\linewidth}{@{}>{\bfseries}p{0.22\linewidth}X@{}}
組織 & - \\
種別 & 論文 \\
時系列 & 2026-05-28 (論文公開) \\
\end{tabularx}
\smallskip
{\bfseries 一次資料から整理したtechnical points}
\begin{itemize}[leftmargin=1.5em,itemsep=0.35em]
\item \textbf{一次情報で確認できる事実}: 保存済み一次資料では「Relevance as a Vulnerability: How Web Retrieval Degrades Safety Alignment in LLM Agents」が2026-05-28に確認できる。
\end{itemize}
\begin{minipage}{\linewidth}
% reader-facing Technical Notes coherent tail group
\begin{itemize}[leftmargin=1.5em,itemsep=0.35em]
\item \textbf{Author claim}: AgentREVEALは、web retrieval integrationと取得contentのrelevanceがLLM agentのsafety alignmentを低下させ得る条件を分析している。
\end{itemize}
{\bfseries 読む際の境界}
\begin{itemize}[leftmargin=1.5em,itemsep=0.35em]
\item \textbf{分析上の留意点}: arXivの原メタデータ／abstractは一次論文資料として扱うが、実験結果と一般化可能性はAuthor claimであり、本号では独立再現していない。
\end{itemize}
\begin{samepage}
{\bfseries 一次資料}
\begingroup\sloppy
\begin{itemize}[leftmargin=1.5em,itemsep=0.25em]
\item {\scriptsize\url{http://arxiv.org/abs/2605.29224v1}}
\end{itemize}
\endgroup
\end{samepage}
\end{minipage}
\endgroup
\end{technicalnote}
\medskip

\begin{technicalnote}{Send a SCOUT First: Pre-hoc Reasoning for Adaptive Detector Allocation in Prompt-Injection Defense}{補足資料}
\begingroup
% reader-facing Technical Notes break policy
\widowpenalty=10000
\clubpenalty=10000
\displaywidowpenalty=10000
\begin{tabularx}{\linewidth}{@{}>{\bfseries}p{0.22\linewidth}X@{}}
組織 & - \\
種別 & 論文 \\
時系列 & 2026-05-29 (論文公開) \\
\end{tabularx}
\smallskip
{\bfseries 一次資料から整理したtechnical points}
\begin{itemize}[leftmargin=1.5em,itemsep=0.35em]
\item \textbf{一次情報で確認できる事実}: 保存済み一次資料では「Send a SCOUT First: Pre-hoc Reasoning for Adaptive Detector Allocation in Prompt-Injection Defense」が2026-05-29に確認できる。
\item \textbf{Author claim}: SCOUTは、requestごとに予測したreliabilityとlatencyに基づき、複数のprompt-injection detectorと必要に応じたLLM judgingを動的に割り当てる方式を提案している。
\end{itemize}
{\bfseries 読む際の境界}
\begin{itemize}[leftmargin=1.5em,itemsep=0.35em]
\item \textbf{分析上の留意点}: arXivの原メタデータ／abstractは一次論文資料として扱うが、実験結果と一般化可能性はAuthor claimであり、本号では独立再現していない。
\end{itemize}
\begin{samepage}
{\bfseries 一次資料}
\begingroup\sloppy
\begin{itemize}[leftmargin=1.5em,itemsep=0.25em]
\item {\scriptsize\url{http://arxiv.org/abs/2605.30837v2}}
\end{itemize}
\endgroup
\end{samepage}
\endgroup
\end{technicalnote}
\medskip
